\subsection{N-012: the printed variable-scaling minimum has no minimum}
\label{sec:nonstandard-viii-scaling}

\paragraph{Precise source statement.}
GCT VIII, subsection 3.3, Conjecture 3.2, printed p.~22
(arXiv:0709.0751v2), varies the rational exponent $a$ in
$c_{a,T}=q^a c'_T$ and minimizes ambient degree over the fibre
\[
 \psi_{a,T}^{-1}(b_{a,T})\subset
 R^{a,\le T}\cap\bar R^{a,\le T}\cap\mathcal B^{\le T}_{\mathbb Q}.
 \tag{N-36}
\]
Here $b_{a,T}$ is the class of $c_{a,T}$ modulo
$q^{1/2}R^{a,\le T}$, and the degree of an ambient coefficient is
its pole order at $q=\infty$ (subsection 3.3, equation (13), p.~21).
The displayed prescription ranges over all rational $a$ and does
not intersect (N-36) with the bar-fixed subspace.  Its surrounding
prose states bar invariance as a goal of globalization; that goal is
not the additional equation $y=\bar y$ in the displayed domain.

Under these literal definitions, the variable-scaling minimum does
not exist.  The defect already occurs on a one-dimensional algebra,
and it occurs on any minimal GT line of the full rank-four algebra.
This is a failure of the displayed prescription.  It does not assert
that no canonical generic basis can be defined by a corrected rule.

\paragraph{The rank-one counterexample.}
Take $\mathcal B=\mathbb Q(q^{1/2})$, $c'_T=1$ and its Laurent
form $\mathbb Q[q^{1/2},q^{-1/2}]$.  Write $A_0,A_\infty$ for
the rings regular at zero and infinity.  For any integer $n\ge0$ set
\[
 a_n=-n,\qquad y_n=q^{-n},\qquad
 R^{a_n}=q^{-n}A_0,\qquad\bar R^{a_n}=q^n A_\infty.
 \tag{N-37}
\]
The element $y_n$ lies in the Laurent form and in $R^{a_n}$;
its quotient by the generator of $R^{a_n}$ is exactly one, so
it maps to $b_{a_n,T}$.  Its quotient by the generator of
$\bar R^{a_n}$ is $q^{-2n}$, which is regular at infinity.
Thus $y_n$ is an admissible candidate in (N-36).  Its unique
nonzero ambient coefficient has degree $-n$.  A nonzero algebraic
function has a finite pole order, so no candidate can have degree
at most $-n$ for every $n$.  Conjecture 3.2 is therefore false
under its displayed unrestricted domain.

This is not an objection based on a scalar-normalization ambiguity:
the initial vector is fixed to be exactly $1$, and the exponents
are ordinary integers already contained in the stated coefficient
field.  Nor does an integral Laurent form remove the example;
every $q^{-n}$ belongs to that form.

\paragraph{The same defect at rank four, for every label order.}
Use the explicit rational GT path lines of N-010, and write
$F=\mathbb Q(D)$, $K=\mathbb Q(v)$ with $q=v^2$.
Let $T$ be any minimal label in the finite partial order used for
globalization.  Then
\[
                  \mathcal B^{\le T}=Kc'_T
 \tag{N-38}
\]
is one-dimensional.  The line is defined over $F$ by N-010, so
it has a nonzero bar-fixed generator $d_T\in\mathcal B_{4,F}$.
Indeed every $P$-word is bar invariant and every scalar in $F$
is fixed by bar.  Clearing the finitely many denominators in an
$F$-word expression gives
\[
 t=g(D)d_T\in\mathbb Q[D]\langle P_1,P_2,P_3\rangle
                  \subseteq\mathcal B_{4,\mathbb Q},
 \qquad \bar t=t,\qquad t\ne0,
 \tag{N-39}
\]
for some nonzero $g(D)\in\mathbb Q[D]$.
The last inclusion uses precisely VIII's ring generated by the
rescaled $P_i$ over its Laurent coefficient form; no saturation
or localization by $g(D)$ is used.

First choose the initial GT vector $c'_T=d_T$ and make $g$ monic,
of degree $m$.  Put
\[
                 a_n=-m-n,\qquad y_n=q^{-n}t
                      \quad(n\ge0).
 \tag{N-40}
\]
Since $D=q+2+q^{-1}$,
$q^m g(D)$ is regular at zero and has value one there.  Hence
\[
 \frac{y_n}{q^{a_n}c'_T}=q^m g(D)\in1+qA_0
                                      \subseteq1+q^{1/2}A_0.
 \tag{N-41}
\]
This proves the required residue condition.  The quotient by
the bar-lattice generator is
\[
 \frac{y_n}{q^{-a_n}c'_T}=q^{-m-2n}g(D),
 \tag{N-42}
\]
whose degree at infinity is $-2n\le0$.  Thus it is regular there.
Every $y_n$ lies in the actual Laurent form by (N-39), so all the
candidates belong to exactly the fibres specified in VIII.

Changing the initial vector on this line does not cure the defect.
For any $c'_T=h(q)d_T$ with $h\in K^\times$, put
$u=t/c'_T$ and $a_0=\operatorname{ord}_0(u)\in\tfrac12\mathbb Z$.
Multiplying $t$ by a nonzero rational constant normalizes the leading
coefficient of $q^{-a_0}u$ to one; this preserves (N-39) and bar
invariance.  For $a_n=a_0-n$ and $y_n=q^{-n}t$,
\[
 \frac{y_n}{q^{a_n}c'_T}=q^{-a_0}u\in1+q^{1/2}A_0,
 \qquad
 \frac{\bar y_n}{q^{a_n}c'_T}=q^{2n-a_0}u\in A_0.
 \tag{N-43}
\]
The second assertion is exactly $y_n\in\bar R^{a_n}$.
Thus the same family is admissible with every initial rational
normalization of the GT line.  Rational powers occurring here are
half-integers, so their powers of $q$ lie in $K=\mathbb Q(v)$.

To finish the argument, express the fixed nonzero $t$ in the ambient
KL tensor basis.  Every nonzero coefficient of $y_n$ has degree
equal to the corresponding degree of $t$ minus $n$, and the support
does not change.  More generally, every nonzero candidate in the
one-dimensional space (N-38) has this same support and differs from
$t$ by one nonzero coefficient-field element of finite valuation.  Its entire finite degree
tuple therefore differs by one finite scalar shift.  Taking $n$
sufficiently large makes the degree tuple of $y_n$ strictly smaller
at every nonzero coordinate.  It is then smaller under VIII's
partial order.  The reverse inequality fails: choose a minimal
nonzero coordinate in the ambient KL order.  Its old degree is
strictly larger, and no preceding coordinate has a smaller old
degree to satisfy the alternative clause in the definition.
The two tuples therefore compare strictly.  No candidate is a
global minimum.  The choice of the partial label order is irrelevant:
every finite nonempty poset has a minimal label, and its line has
the property just proved.

\paragraph{What this decides, and what a repair would change.}
The literal external minimum-degree prescription of VIII does not
construct $E(4)$, even after all GT matrices, path labels and initial
normalizations have been supplied.  Consequently there is no full
table singled out by that prescription whose signed constants,
coefficient field or dimension ring can be computed as requested.
The displayed rank-three examples and the one distinguished
rank-four element remain explicit elements and can be studied on
their own; this argument does not invalidate their algebraic formulas.

One possible repair is to require $y=\bar y$ explicitly.  That
excludes $q^{-n}t$ for $n>0$ because $\bar t=t\ne0$.  Another is
to bound or fix the permitted scaling exponents.  Neither repair
is part of (N-36), and neither is proved here to yield unique
minimum lifts or the printed examples.  VIII's section 4 gives a
different internal degree notion and explicitly begins its
discussion by assuming a bar-invariant element; our counterexample
does not refute that additionally restricted alternative.  The
text conjectures agreement of the two procedures, which cannot be
taken as established in the presence of the displayed external
defect.

Thus item 9(a) has a precise obstruction in its stated source
prescription, rather than a missing numerical computation.  A
repaired canonical-basis problem remains distinct from both the
proved integral nonfreeness in N-009 and the proved rational GT
construction in N-010.

\paragraph{Checks and literature boundary.}
The source passages checked are VIII's coefficient rings in
subsection 2.2, the whole of subsections 3.1--3.3, equation (13),
Conjectures 3.1 and 3.2, and the separate internal definition in
section 4.  The proof above is an elementary algebraic
counterexample, independent of a numerical experiment.  Its
rank-four extension uses N-010's independently reviewed rational
GT construction; the rank-one counterexample has no such input.
Targeted searches for the exact title with ``minimum'' and
``counterexample'', for ``Mulmuley minimum degree scaling canonical
basis'', and for ``0709.0751 Conjecture 3.2'' found no earlier report
of this defect.  The appropriate historical claim is limited to that
search and the supplied sources.  The plethysm branch independently
read the whole proof and the complete primary source passages;
its review is recorded in
\nolinkurl{nonstandard_plethysm_review_N012.md}.
